\documentclass[10pt,letterpaper]{article}
% Add a bunch of useful math, font, and symbols
\usepackage{amsfonts}
\usepackage{amsmath}
\usepackage{amssymb}

% English support
\usepackage[english]{babel}

% Line break for long URLs
\usepackage{breakurl}

% Citation
\usepackage[superscript]{cite}

% Better enumerate and itemize
\usepackage{enumitem}

% Better control of headers and footers
\usepackage{fancyhdr}

% Floating objects like figures and tables
\usepackage{float}

% Page layout and dimensioning
\usepackage[margin=1in]{geometry}

% Basic color, graphics and text manipulation 
\usepackage{graphicx}

% Use Helvetica typeface
\usepackage[scaled]{helvet}
\renewcommand\familydefault{\sfdefault} 
\usepackage[T1]{fontenc}

% Cross-referencing hyperlinks
\usepackage{hyperref}

% Accept utf-8 input encoding
\usepackage[utf8]{inputenc}

% Make indexes
\usepackage{makeidx}

% Microtype (apparently makes the typographics stuff better)
\usepackage{microtype}

% [Disabled] multi-column writing
% \usepackage{multicol}

% English ordinal counting (1st, 2nd, etc.)
\usepackage{nth}

% Long table
\usepackage{longtable}

% Paragraph skip - adds extra lineskip spacing
\usepackage{parskip}
\setlength{\parskip}{0.7\baselineskip plus 2pt}

% Add ability to set space between lines
\usepackage{setspace}

% S.I. units
\usepackage{siunitx}

% Subcaptions for subfigures
\usepackage{subcaption}

% Include svg graphics
\usepackage{svg}

% Drawing graphics
\usepackage{tikz}

% Subsubsubsection
\usepackage{titlesec}
\setcounter{secnumdepth}{4}
\titleformat{\paragraph}
{\normalfont\normalsize\bfseries}{\theparagraph}{1em}{}
\titlespacing*{\paragraph}
{0pt}{3.25ex plus 1ex minus .2ex}{1.5ex plus .2ex}

% Custom titles
\usepackage{titling}

% Url
\usepackage{url}

% Custom definitions
\newcommand{\doctitle}{Design Specification}
\newcommand{\docsubtitle}{Computing Platform Multirotor with FPGA Hardware Acceleration Applications}

% Custom commands
\newcommand{\ts}{\textsubscript}	% Subscript command %
\newcommand{\revisionnum}{!! Revision number not found. Please add renew command \\revisionnum in the document after the ``aliases.tex'' include !!}
\renewcommand{\revisionnum}{1.3}

% Use hyphans to break up urls
\def\UrlBreaks{\do\/\do-}

% PDF and href setup
% Hyper ref
\hypersetup{
	colorlinks=true,
	citecolor=black,
	linkcolor=black,
	filecolor=black,
	urlcolor=blue,
	pdftitle={\@title},
	bookmarks=true
}
\urlstyle{same}
% Page headings
\pagestyle{fancy}
\fancyhead[L]{\MakeUppercase{CPEN/ELEC 491}}
\fancyhead[C]{\textbf{\doctitle}}
\fancyhead[R]{Mieszko Lis, PhD}
\fancyfoot{}
\fancyfoot[C]{Non-Confidential}
\fancyfoot[R]{\thepage}

% No paragraph indent
\parindent 0ex

% Meta
\author{
	Deutsch, Peter &
	\textit{me@peterdeutsch.ca}
	\\
	He, Muchen &
	\textit{i@muchen.ca}
	\\
	Hsueh, Arthur &
	\textit{ah11962@outlook.com}
	\\
	Wang, Meng &
	\textit{wzfftxwd@gmail.com}
	\\
	Wilson, Ardell &
	\textit{ardellw96@gmail.com}
}
\title{\doctitle}
\date{\today}
\makeatletter
\renewcommand{\maketitle}{
	\bgroup
	\setlength{\parindent}{0pt}
	\begin{flushleft}
		% Top spacing
		\vspace*{0.5in}

		% Team logo
		\includegraphics[scale=0.4]{../assets/capstonelogo1.png}
		\vspace*{0.25in}

		% Title
		\textbf{\Huge{\@title}}\\
		\hrulefill

		% Subtitle
		\textbf{\huge{\docsubtitle}}
		
		\vspace*{0.5in}

		% Course number and team
		\textbf{\Large{CPEN/ELEC 491 Capstone Team 109}}\\
		\hspace*{0.1cm}
		\begin{tabular}[h]{|ll}
			\@author
		\end{tabular}

		\vspace*{0.25in}

		{\Large{\textbf{Dr. Mieszko Lis}, Client}}\\
		\hspace*{0.1cm}
		\begin{tabular}[h]{|ll}
			Electrical and Computer Engineering, The University of British Columbia
		\end{tabular}

		\vfill
		
		% Date
		\large{Revision 1.0 -- \@date}
		\vspace*{0.5in}

		% Logo
		\hspace*{-0.3cm}\includegraphics[scale=0.5]{../assets/ece_logo.pdf}

	\end{flushleft}
	\egroup
}
\makeatother

% Begin Document
\begin{document}
% Title Page
\begin{titlepage}
	\maketitle
\end{titlepage}


\renewcommand{\thepage}{\roman{page}}
\setcounter{page}{1}

% Revision history
\backgroundsetup{
	scale=1,
	color=black,
	opacity=0.2,
	angle=0,
	contents={\includegraphics[height=\paperheight,width=\paperwidth]{../assets/render1bg}}
}
\BgThispage
\thispagestyle{empty}
\section*{Revision History}

Revision history written here.

\begin{table}[H]
\begin{tabular}{llll}
\hline
Version \# & Initials & Release Date & Changes Made \\ \hline
0.0 & PD & 2019-10-11 & Initial skeleton of the document.\\
\hline
\end{tabular}
\end{table}
\clearpage

% Table of contents
\setcounter{secnumdepth}{3}
\tableofcontents
\thispagestyle{empty}

% Roman numeral page numbers


% Terms and Abbreviations
\thispagestyle{empty}

\section*{Terms and Abbreviations}

\begin{tabular}[h]{rp{0.75\linewidth}}
    \hline
    \textbf{Term} & \textbf{Definition}\\
    \hline

    ANN & Artificial Neural Network, or simply ``Neural Network'', is a data processing model modeled after neuron interactions. The process consists of forward propagation through the use of several matrix multiplications.\cite{ann}\\
    ASIC & Application-Specific Integrated Circuit.\\
    CNN & Convolutional Neural Networks are neural networks which are especially useful for image classification.\cite{cnn} \\
    ECE & Department of Electrical and Computer Engineering at the University of British Columbia.\\
    FPGA & Field-Programmable-Gate-Array, a ``programmable'' hardware unit that allows for ASIC-like performance with software-like turn-around time and flexibility.\\
    FPS & Frames Per Second.\\
    GPU & Graphics Processing Unit, a discrete piece of hardware designed to accelerate graphics-intensive or other parallel computing tasks.\\
	ICICS & Institute for Computing, Information, and Cognitive Systems at the University of British Columbia.\\
    LOS & Line-of-sight.\\
    ML & Machine Learning.\\
    Multirotor & An unmanned vehicle with multiple engines. \\
    OTS & Off-the-shelf, or commercially available/purchasable. \\
    PID / PID Controller & Proportion-Integral-Derivative controllers denote the most common control algorithm for precise and accurate movement in multirotor applications.\cite{pid}\\
    RNN & Recurrent Neural Networks are neural networks where the output depends on previous computations, effectively consisting of memory.\cite{rnn}\\
    RX & Receiver.\\
    TC & Transport Canada.\\
    TX & Transmitter.\\
    YOLO & You-Only-Look-Once is a fast ML algorithm which detect objects, but is unlike CNN or RNN.\cite{yolo}\cite{yolo-2}\\
     & \\

    \hline

\end{tabular}

Technical terms and abbreviations dictionary go here.


% List of figures and tables
\iftotalfigures
\addcontentsline{toc}{section}{\listfigurename}
\listoffigures
\fi
\iftotaltables
\addcontentsline{toc}{section}{\listtablename}
\listoftables
\fi

\newpage

% Set page and section counter
\renewcommand{\thepage}{\arabic{page}}
\setcounter{page}{1}

% TODO: fill out all the sections
% if the sections gets too long, move them to a separte .tex document
\section{About This Document}
\subsection{Purpose}
This document serves to provide a concise description of the primary design decisions made throughout the Multirotor Computing Platform project. In addition to providing a high-level overview of the system architecture, this design specification provides specific rationale for the design of each of the project's various subsystems. This document also outlines potential improvements to subsystems which were envisioned, but not implemented due to the COVID-19 pandemic.

\subsection{Intended Audience}
This document is intended for the:
\begin{itemize}
\item Project Designers and Quality Controllers (the Capstone team),
\item Client (and their representatives),
\item External Project Sponsors,
\item and Legal Personnel (The Department of Electrical and Computer Engineering, Industry Canada, Transport Canada)
\end{itemize}

This document is particularly relevant to those who may modify, extend, improve, or interface with the system, in addition to those who will assess the product's design against academic, organizational, and/or legal criteria.

\subsection{Reading Guide}

This document begins with a high-level decomposition of the product into its two major components: the \textit{Computing Platform} and the \textit{Multirotor}. It then conveys the technical decisions made for the subsystems of each component, providing detailed rationale as to why they represent the best solution given the project's requirements and constraints (as detailed in the \textit{Requirements Document}). Cross-references to the \textit{Requirements Document} are denoted in \textbf{bold}.

It is useful to note that the \textit{Computing Platform} subsections (beginning in Section \ref{computing_platform}) are ordered by their stage in the device's data flow -- video capture is discussed first, then processing, and finally data conveyance to the user.

\newpage
\section{High Level Design}\label{high_level_design}
The Computing Platform, mounted on the multirotor, is comprised of two distinct processing platforms: the \textit{Platform Management Board} (PMB) and the \textit{Programmable Logic Board} (PLB). The PMB hosts the software and interfaces required to interface with both the PMB (also present on the multirotor) and the base station.

\begin{figure}[H]\label{hlpic}
    \centering
    \includegraphics[width=\linewidth]{img/intpic}
\caption{High-Level System Integration with RPAS}
\end{figure}

\begin{figure}[H]\label{hldiag}
\centering
\includegraphics[width=15cm]{img/highlevel.pdf}
\caption{High-Level System Architecture}
\end{figure}



\section{Computing Platform Technical Design}\label{computing_platform}

\subsection{Video Capture}\label{video_capture}
% TODO: Talk about camera choice, MIPI interface, camera libraries. Also, talk about alternatives we explored (i.e. analog video transmission/acquisition.

The video capture subsystem of the computing platform consists of a camera, its video transmission protocol to the processing core and video transmission to the base station. For this subsystem, the hardware used includes a CMOS camera, a Raspberry Pi device, and MIPI Camera Serial Interface protocol. Design description of the Raspberry Pi device is described in the next section.

\subsubsection{Camera Choice}
The camera used to capture video is the Raspberry Pi Camera Module. It supports 1080p, 720p and 640x480p video. It has a fixed focus lens. The camera will capture live video and send the video data to the Raspberry Pi, which will feed the data in different forms to both the machine learning model and the base station.

% camera Choice
The decision of the camera was influenced by the requirement of future replaceability. The project's goal is to create a system which provides both video acquisition and supplemental programmable hardware processing capabilities. The Raspberry Pi device contains different ports which support different cameras with varying characteristics and performance.

One major consideration regarding selecting a camera is the latency and throughput of the machine learning implementation on the FPGA. As the machine learning implementation will be heavily constrained by the FPGA's limited resources, it is will not be necessary to use a camera with and especially resolution and framerate.

% {MIPI interface}
\subsubsection{MIPI Interface}
The MIPI camera serial interface (CSI) is a high-speed protocol primarily intended for point-to-point image and video transmission between cameras and host devices. 

%requirement got removed ...
The MIPI protocol is a widely adopted protocol for image and video transmission, fullfilling requirement (xx). A hardware based MIPI core will be programmed into the programmable logic and will handle receiving of camera footage.

% {Camera libraries or getting the video to a 'viewable' state}
\subsubsection{Camera Libraries}
Live video feed is a requirement for the base station. This is achieved by implementing a Python script on the Raspberry Pi for streaming the video to a webpage on the base station via Wi-FI. Justification for using Wi-FI is discussed in section~\ref{sec:Communications}. The script used to achieve this uses the Python package 'picamera,' which is a Python interface to the Raspberry Pi Camera. The 'picamera' package is a widely supported open-source resource fullfilling requirement (xx). It is also well documented, so the client will be able to easily adjust interfacing of the camera to their needs.

% {Explored Alternatives}
\subsubsection{Explored Alternatives}
There were many alternatives explored for each component of the Video Capture subsystem.

One of the possibilites for the camera was purhcasing a drone package with a pre-integrated camera. The implicit requirement to disassemble and reconfigure/tap-into a hard-wired system, however, would create an unmaintainable product -- not allowing for user replacement of the camera for future research improvements (violating constraint XX). %again the requirement got removed

Another possible choice for the camera was purchasing a USB webcam and connecting it to the USB port to the Raspberry Pi device. The issue with this option is the inferior configurability of the webcam compared to the Raspberry Pi camera. The reconfigurability provided by the 'picamera' Python package, as previously mentioned, made the Rasberry Pi Camera a better solution suited to the client's needs.

For the live video feed transmission, other options explored include wireless HDMI, VGA and analog transmission methods. Again, detailed justification for using Wi-FI and not these is discussed in a later section but the decision came down to a few factors. The original options required additional hardware (higher costs and system weight) and would require more complexity for the user to handle.


\subsection{Platform Management Board}\label{platform_management}
TODO: Why the RPi? Other boards considered? What kind of video encoding can we do? Why are we using WiFi? Justify why we have two boards.

\subsection{Programmable Logic Board}\label{programmable_logic}
\subsubsection{Subsystem Overview}
The Programmable Logic board manages the hardware-accelerated ML application, performs additional pre/post-processing on ML data (if required by the ML application), and communicates with the Platform Management board.

A decomposition of the system is seen on the right in Figure \ref{pcdiag}.

Of note is the distinction between the two portions on the Programmable Logic board: the Hard Processor System (HPS) and the FPGA. The Hard Processor System is a non-customizable CPU and is used for generic tasks (such as managing I/O via the Ethernet port and processing non-acceleratable ML tasks). The FPGA provides the system's \textit{programmable circuity} -- performing the requisite hardware acceleration needed for ML tasks. The two systems co-reside on the same chip, and are interfaced via an internal bus.

\subsubsection{Device Selection}
The selected Programmable Logic board is the \textit{ZedBoard SoC Development board} (pictured in Figure \ref{zedboard}), manufactured by Digilent. 

\begin{figure}\label{zedboard}
\centering
\includegraphics[width=12.5cm]{img/zedboard_functional_overview.jpg}
\caption{Layout of Connectors and Main ICs on ZedBoard}
\end{figure}

The system specifications are as follows:
\begin{itemize}
\item Xilinx Zynq-7000 AP SoC XC7Z020-CLG484
\item Dual-core ARM Cortex-A9 (\textit{HPS})
\item 512 MB DDR3 Random Access Memory
\item 256 MB Quad-SPI Flash Memory
\item 10/100/1000 Ethernet 
\item USB OTG 2.0 and USB-UART 
\item PS \& PL I/O expansion (FMC, Pmod, XADC)
\end{itemize}

The selection of this particular Programmable Logic board was mainly motivated by cost. The client was able to provide several ZedBoard units free-of-charge, resulting from donations through the Xilinx University Program. While this board has a fairly small number of logical cells (85,000 LCs) compared to other programmable logic boards on the market, especially those which target video applications, the project's budget does not allow for the purchase of more performant boards without significant sacrifices to the multirotor assembly. The number of logical cells on the ZedBoard is sufficient to comply with the minimum LC constraint (\textbf{C.EX.1}).

Examples of other boards examined for this project include the Nexys Artix 7 FPGA (200,000 LCs, US\$479) and the Microsemi PolarFire Video/Imaging FPGA (300,000 LCs, CAN\$1500). Use of the Terasic DE1-SoC (another board available free-of-charge from the client, 85,000 LCs) was also considered, but was ultimately not selected due to its poorer I/O selection (no FSC/HDMI connections, which could be used by the client to expand the system) and difficult-to-use development toolchain.

\subsubsection{Data Flow}
The overall data flow managed by the PLB is as follows:
\begin{enumerate}
\item Video/ML input data is acquired by the HPS (ARM core) via the Ethernet port (from the PCB)
\item The HPS, managing the machine learning, splits the input data into ML \textit{subtasks} (such as matrix-matrix multiply, pooling, etc)
\item Where hardware acceleration of a particular ML subtask is possible, the requisite data is sent from the HPS to the FPGA via a bi-directional AXI interface
\item The ML IP block on the FPGA receives subtasks to be accelerated via the AXI bus, performs the requisite processing, and sends the results to the HPS (on the same bus)
\item When all ML subtasks are complete, the HPS sends the results of the task back to the PCB (Raspberry Pi) via the Ethernet port
\end{enumerate}

\subsubsection{ML Hardware Accelerator Implementation}
To facilitate the hardware acceleration of ML subtasks, a customized version of the open-source MARLANN (\textit{Multiply-Accumulate and Rectified-Linear Accelerator for Neural Networks}) module will initially be used. As the client intends to replace this module with a custom-designed solution stemming from their research, however, it is only intended as a placeholder.

MARLANN is capable of performing 8-bit fixed point signed integer arithmetic, matrix multiplication, max-pooling, and applying the ReLU activation function. After evaluating other open-source FPGA-based accelerators, MARLANN was selected due to its low FPGA area overhead (approx. 5,000 LCs), efficient pipelined architecture, and comprehensive documentation. The use of a pre-designed accelerator also allows the developers to quickly deploy simple ML models without the considerable hardware design overhead implicit in a full-custom solution.

\subsection{Machine Learning Model}\label{machine_learning}
% TODO: Describe potential machine learning applications. What could this look like, what will this \textit{not} look like? What kind of accelerators already exist? What might our solution look like?

% 


Machine learning is the topic that provides the context of the project. This project is not focused on creating a machine learning model, but creating a mobile computing platform that will allow the client to perform research in machine-learning models.

A mobile computing platform means we are able to perform processing in unaccessilble locations. Using a drone, we are creating an unmanned aerial vehicle with the ability for its user to change operating functions using programmable logic.

A user can program the logic to perform one application, allow the drone to perform that function autonomously, and then program the logic for a completely differnent application, without having switch its processing hardware.

A computing platform with machine learning means that it can automatically learn and improve from experiences without being explicitly programmed. This concept is accentuated when the computing platform becomes mobile; the system can autonomously serve its purpose without needing physically close monitoring.

\subsubsection{Correct Machine Learning Applications}
Because this is a mobile platform, the computing platform serves to take in physical data from the external world, process it, and output data that can be analyzed. Examples of this usage would include image/video recognition, as well as pattern recognition and classification in any environment.

\subsubsection{Incorrect Machine Learning Applications}
As a result, the computing platform would not serve its purpose in applications where the data to be processed is sent its user. The mobile computing platform should not be used in applications such as email classification, spam filtering and virtual assistants.

Also factoring in the mobility of the system, the machine learning implementation should not be used in applications should be performed on the ground (even if ml can be applied.) Cases of grounded applications include medical services and speech recognition.

\subsubsection{Hardware Accelerators} %????
The hardware accelerators used for machine learning can be applied in CPUs, GPUs, FPGAs and ASICs. 

\subsubsection{Project Implementation}
The machine learning implementation is not the focus of this project but it is required to demonstrate the capabilities of the mobile computing platform. 

Our solution would be a simple hardware accelerated machine learning algorithm capable of processing live and test video and outputting metadata that can be analyzed. This can come in the form of simple object detection. The deep learning approach for this application is the use of Convolutional Neural Networks.



\subsection{Communications}\label{communications}
\subsubsection{Subsystem Overview}
The Air-Ground communication scheme between the platform management board and the base station is responsible for transmitting video, ML results, and control commands. It is a bi-directional, high bandwidth communication channel that requires a stable connection range up to 30 meters away (\textbf{NF.CM.1}).

\subsubsection{Single-Link versus Dual-Link Solution}
There are two options for establishing this communication link - single channel or dual channel. 

The first option is to use the existing communication channel to the multirotor (used for drone control), to transmit data. This would be advantageous as it doesn't require additional hardware and saves power. The disadvantage is that the drone control is using an off-the-shelf module, and in order to extend its functionality we would need to make modifications to a highly integrated system. 

The second option is to use a separate air-ground communication channel. This channel can be established under commonly used wireless communication methods (ex. Wi-Fi or Bluetooth), requiring a dedicated pair of transceivers. The advantage of this option is that most development boards are equipped with wireless transceivers and the communication scheme is easily configured. As such, we decided to employ this option.

\subsubsection{Bandwidth Requirements}

The majority of the bandwidth will be reserved for video transmission as it's the most intensive transmission task performed by the computing platform. The total number of bits per second of video that need to be transmitted can be calculated using the following formula:

$$
\text{\# of bits/second} = \frac{\text{\# of bits}}{\text{1 pixel}} \times \frac{\text{\# of pixels}}{\text{1 frame}} \times \frac{\text{\# of frames}}{\text{1 second}}
$$

The video that will be transmitted, assuming no compression, has a resolution of 640x480 pixels at 30 frames per second. The colour depth is 24-bit for a RGB colour set (that is, each pixel is 8 bits). As such:

$$
\text{\# of bits/second} = \frac{8}{\text{1 pixel}} \times \frac{640\times 480}{\text{1 frame}} \times \frac{30}{\text{1 second}} = 73.738~\text{Mbits/s}
$$

The other data that needs to be transmitted includes the video metadata, the results from ML model, the control commands, and other protocol overheads. Thus, an upper bound for the total bandwidth required by the system is roughly 75 Mbps.

\subsubsection{Communication Scheme}

The communication method we chose for the project is Wi-Fi. The specific standard we chose to utilize is IEEE 802.11n (Wi-Fi 4), as it has enough throughput and nearly all modern devices equipped with Wi-Fi support this standard. (\textbf{C.EX.3})

The advantages of Wi-Fi over other communication methods (such as Bluetooth) are as follows:
\begin{itemize}
    \item Long transmission range. IEEE 802.11n Wi-Fi has a transmission range of up to 53 m. (\textbf{NF.CM.1})
    \item High transmission rate. IEEE 802.11n Wi-Fi has a transmission rate up to 600 Mbps.
    \item No license is required.
    \item After the initial configuration, the connection is capable of auto-establishment. (\textbf{F.CM.3})
\end{itemize}

\subsection{Base Station}\label{base_station}
\subsubsection{Subsystem Overview}

The base station is a mobile device that connects to the computing platform on the multirotor. In particular, it performs the following tasks:
\begin{itemize}
    \item Displays the status of the system
    \item Displays the ML results in real-time (by overlaying the results over live video)
    \item Stores the log data and video for further analysis
    \item Remotely configures the computing platform (ex. starting object detection)
\end{itemize}

\subsubsection{Device Alternatives}
We considered two classes of devices which could be used as a base station: laptops and smart mobile devices (such as cellphones or tablets). The common features of these two options are: they are both equipped with a screen which can be used to display the video and ML results, they both have a Wi-Fi transceiver (\textbf{F.BS.3}), and they both have non-volatile storage which can be used to store results (\textbf{F.BS.6}).

One functionality that mobile devices lack is the support for common machine learning platforms (ex. TensorFlow). As the client may want to perform additional ML processing on the base station for further analysis in the future, laptops were ultimately selected as the base station's platform. Support for smart devices (to facilitate multiple viewers) may be added if time permits.

\subsubsection{Functionality}
A web-application is used as the GUI for the base station. This allows for universal compatibility with all laptops, regardless of their operating system.

The base station hosts a TCP/IP server, and serves two clients: the PMB and the web-application. The server continuously queries the PMB for live video data and ML metadata results, and sends the data to the web-application for display.

The ML results are displayed as a bounding box around the detected object(s) in the frame (\textbf{F.BS.1}). In particular, the ML results consist of three parts: the class of the detected object(s), the coordinates of the upper-left vertex of the bounding box(es), and the height/width of the bounding box(es).

The web application also displays real-time system status messages. These messages come from the PMB or the base station itself. The status updates include information on initiation or completion of processes, changes in PMB, PLB or base station settings, and system warnings/errors.

The base station continuously saves the received video and metadata locally for any analysis purposes (\textbf{F.BS.6}). All faults are also logged locally on the device to facilitate further debugging (\textbf{F.BS.5}).
%% maybe we just log all status messages, not limited to warnings/errors.

The user may use these buttons and sliders on the web-application to adjust the visual properties of the web-application or change the settings on the PMB and PLB (\textbf{F.BS.2, F.BS.4}) such as the camera polling rate and ML sensitivity.

\subsubsection{Programming Libraries}

There are two JavaScript libraries used in both the client-side and server-side code of the base station.

The first library is the socket.io library, which enables real-time, bidirectional and event-based communication between browser and server. This library is used to ensure a single communication protocol between the base station server and its clients. The client-side portion of the library is available in other languages such as C++ and Python with the same communication protocol, allowing for more versatility when programming the PMB to communicate with the base station. 

The second library is the p5.js library, which is a client-side JavaScript platform that provides programmers with tools to create virtual canvases more easily. This library is used to handle live video feed and ML bounding box data display; it allows for an elegant and high-level approach to programming this part of the web application. 


\subsection{Computing Power}\label{computing_power}
\subsubsection{Subsystem Overview}
The computing platform's power subsystem consists of a commercially produced integrated battery and power management pack. The devices powered by this subsystem include the:

\begin{itemize}
\item Platform Management Board (Raspberry Pi) -- maximum current draw of 0.5 A and a typical current draw of 0.4 A, and
\item Programmable Logic Board (Zedboard) -- maximum current draw of 1 A and a typical current draw of 0.4 A
\end{itemize}

The power system must provide continuous power to the computing platform for at least 20 minutes (\textbf{F.PR.1}).

\subsubsection{Seperate vs Common Power Supplies}
Two power supply schemes were considered for the device: a common power supply shared between the multirotor and the computing platform, or separate power supplies across the two systems. While separating the power systems adds complexity and weight, it provides the following benefits:

\begin{itemize}
    \item No risk of noise from the motors interfering with the computational hardware.
    \item Allows for testing of the computational hardware separately from the drone.
    \item Reduces the difficulty of replacing the multirotor assembly, if required.
\end{itemize}

Ultimately it was determined that the advantages of using separate power supplies, in particular, avoiding motor noise, warranted the extra weight and complexity.

\subsubsection{Custom vs Commercial Power Supply}
Initially a custom designed power supply system was intended to be used. It was to consist of a lithium-polymer or lithium-ion battery and a battery management system. The battery management system (seen in Figure \ref{powerdiag}) would have consisted of two DC-DC converters (with an approximate efficiency of 94\%), a battery protection board, and a battery monitor. 

\begin{figure}[H]
\centering
\includegraphics[width=15cm]{img/Power_Diagram.png}
\caption{Original Custom-Designed Computing Power System}
\label{powerdiag}
\end{figure}

DC-DC conversion would have been accomplished through use of a buck converter for the 5 V output and a buck-boost converter for the 12 V output. 

After designing the system, an alternative commercial system was found which fits the project requirements. Said commercial system consists of 3 lithium-ion batteries with a combined capacity of 3000mAh in addition to a power management circuit. The power management circuit provides a 5V USB output and a 12V output which are compatible with the Raspberry Pi and the Zedboard respectively. The rated current output of 3A is sufficient to simultaneously supply both computation boards, and the battery capacity is rated to provide 3 hours and 45 minutes of operation on a single charge.

Ultimately, the commercial solution was chosen over the custom-designed solution as the commercial system provides all of the requisite functionality at a lower cost and smaller form factor.


\subsection{Multirotor RPAS}\label{multirotor}
\input{3_2_multirotor_V2}

% Bibliography
\clearpage
\addcontentsline{toc}{section}{References}
\bibliographystyle{ieeetr}
\bibliography{references}

% Appendix (uncomment to enable appendix)
\clearpage
\appendix
\section{Multirotor Design}\label{appendix:multirotor}
% \subsection{Multirotor Configurations}
% Since the content in the multirotor section is large and generous amount of information is given as to not assume the knowledge of the reader. Some of the text can be put here instead.

\subsection{Multirotor Motor Speed and Size}\label{appendix:droneengine}

All BLDC outrunner motors have the same height, and on product pages they are often given the code 22XX to denote their stator size. The design parameter is the radius (XX on the product code). Typically, the larger the stator (radius), the slower it spins given the same voltage. This is a linear relationship described by the motor parameter kV:

$$
\omega_{\mathrm{RPM}} = \mathrm{kV} \times \mathrm{Voltage}
$$

The kV constant is inversely proportional to another motor constant, kT, which models the linear relationship between current and torque:

$$
\tau = \mathrm{kT} \times \mathrm{Current}
$$

The inverse relationship can be shown by equating electrical power into the motor and mechanical power produced by the motor:

$$
\tau\omega_{\mathrm{RPM}} = \mathrm{kV}\mathrm{kT}\times VI
$$

By constraining the electrical power, we can make trade-off between torque and rotation speed.
Air resistance (“drag”) is proportional to speed of an object squared, therefore the power required to 
lift increases non-linearly to the force output. As seen in the test data obtained from one manufacturer’s 
datasheet below, when using the same propellers we see an increased thrust output resulting from increased voltage (and 
thus current). However the efficiency, typically measured in g/W, decreases. 

Below is a sample set of thrust output using a SunnySky X2216 1250 kV motor using 9.0 inch 5.0 inch pitch and 10.0 inch 4.7 inch pitch propellers obtained from a vendor website \cite{sunnysky-2216}. 

\begin{table}[H]
    \centering
    \caption{Sample Motor Test Data}
    \label{table:sunnyskyx2216-table}

    \begin{tabular}{lrrrrr}

    \hline
    \textbf{Propeller} & \textbf{Voltage} & \textbf{Current} & \textbf{Pull}  & \textbf{Power Consumption} & \textbf{Efficiency}\\
    & [V] & [A] & [g] & [W] & [g/W] \\
    \hline
     & 10.0 & 16.4 & 940 & 164 & 5.73 \\
    9050 & 11.0 & 19.5 & 1050 & 215 & 4.89 \\
     & 12.0 & 21.5 & 1250 & 258 & 4.84 \\
    \hline
     & 10.0 & 23.5 & 1200 & 235 & 5.10 \\
    1047 & 11.0 & 26.5 & 1350 & 372 & 3.63 \\
     & 12.0 & 30.2 & 1520 & 363 & 4.19 \\
    \hline

    \end{tabular} 
\end{table}

Therefore, it makes sense that we choose larger motors with relatively lower kV constants. Slower motors 
are also safer since there is less angular momentum of the blades. Motors with 800 kV to 1300 kV are 
suitable for our applications and achieve reasonable thrust output with adequate efficiency. 2212, 2214, 
and 2216 stator sizes are suitable in this range of kV.

\end{document}